\documentclass{article}
\usepackage[T1]{fontenc}
\usepackage{lmodern}
\usepackage{graphicx}
\usepackage{amsmath,amssymb}
\usepackage[a4paper,margin=2cm]{geometry}
\usepackage[absolute,overlay]{textpos}
\setlength{\TPHorizModule}{1mm}
\setlength{\TPVertModule}{1mm}
\usepackage[indonesian]{babel}
\usepackage{hyperref}
\setlength{\emergencystretch}{2em}
% unit: Halaman 17

\begin{document}

% === Halaman 17 (llm) ===

\begin{itemize}
    \item Diperoleh pemetaan (cipherteks $\rightarrow$ plainteks):
    \[
    \begin{aligned}
    P \rightarrow e \\ 
    Z \rightarrow t \\ 
    W \rightarrow h \\ 
    S \rightarrow a
    \end{aligned}
    \]
    \item Iterasi 2:
\end{itemize}

\vspace{5mm}

\begin{verbatim}
UZ QSO VUOHXMOPV GPOZPEVSG ZWSZ OPFPESX UDBMETSX AIZ
t a e e te a that e e a a t

VUEPHZ HMDZSHZO WSFP APPD TSVP QUZW YMXUZUHSX\ne t ta t ha e ee a e th t a

EPYEP OPDZSZUFPO MB ZWP FUPZ HMDJ UD TMOHMQ\ne e e tat e the et
\end{verbatim}

\end{document}
